\documentclass[11pt,a4paper]{article}
\usepackage[utf8]{inputenc}
\usepackage[margin=0.8in]{geometry}
\usepackage{amsmath,amssymb,amsfonts}
\usepackage{booktabs}
\usepackage{graphicx}
\usepackage{hyperref}
\usepackage{caption}
\usepackage{microtype}
\usepackage{xcolor}
\usepackage{placeins} % Force tables to stay within sections

\hypersetup{
    colorlinks=true,
    linkcolor=blue!70!black,
    citecolor=blue!70!black,
    urlcolor=blue!70!black
}

\title{\textbf{Integration Order, Cointegration, and the Specification Fork for Mechanization Intensity: A Unified Empirical Audit for the United States (1925--2021) and Chile (1920--2024)}}
\author{\textbf{Chapter 2 Research Group} \\ Department of Economics, UMass Amherst}
\date{July 2026}

\begin{document}

\maketitle

\begin{abstract}
This paper examines the integration orders and cointegration properties of heterogeneous capital---machinery ($K^{\text{ME}}$) and nonresidential structures ($K^{\text{NR}}$/$K^{\text{NRC}}$)---across century-long samples in the United States (1925--2021) and Chile (1920--2024). We combine non-parametric LOESS 2nd derivative inflection dating, the complete R \texttt{urca} unit root battery on levels and first differences, Johansen VECMs under drift and trend stationarity controls, i.i.d.\ residual bootstrap resampling, and Cavaliere et al.\ (2012) wild bootstrap rank testing for non-stationary volatility robustness. Cointegrating vectors are normalized strictly with Nonresidential Structures ($k^{\text{NRC}}$/$k^{\text{NR}}$) on the left-hand side and Machinery ($k^{\text{ME}}$) on the right-hand side ($k^{\text{NRC}} = \alpha + \gamma k^{\text{ME}}$). Three empirical results stand out. First, log-levels of composing capital assets are non-stationary $I(1)$ processes; first differences are strictly stationary. Second, under drift controls, un-augmented century VECMs yield cointegrating coefficients near unity ($\hat{\gamma} = 0.994$ for the U.S.; $\hat{\gamma} = 0.955$ for Chile), encoding balanced Fordist accumulation where structures tracked machinery expansion proportionally. Third, speed-of-adjustment coefficients ($\alpha$) reveal structural asymmetry: in the U.S., both assets adjust to disequilibrium, with machinery adjusting twice as fast during Fordism ($\alpha_{\text{ME}} = -0.323$ vs.\ $\alpha_{\text{NRC}} = -0.151$); in Chile, machinery is weakly exogenous ($\alpha_{\text{ME}} = -0.011$, $p > 0.40$), forcing structures to absorb the full burden of adjustment to foreign-exchange-constrained technique choice.

Building on this evidence, we provide a definitive research verdict on the specification of mechanization intensity for the $\theta_t$ estimation stage. Six plausible candidates---unrestricted component levels, the log-ratio $\tau_t$, raw mechanization intensity $q_t = K^{\text{ME}}/L$, and three accumulated-path proxies ($Q_t^\omega$, $Q_t^{\text{ME share}}$, $Q_t^q$)---are evaluated against theoretical admissibility from the choice-of-technique micro-meso framework and integration order evidence. Only the capital composition state $\tau_t = \ln K^{\text{ME}} - \ln K^{\text{NRC}}$---entering through residual-centered specification $D_9$ alongside the nonresidential-capital envelope---survives both gates. Unrestricted component levels, raw $q_t$, and accumulated-path proxies are rejected on econometric grounds; their theoretical motivation does not override stochastic non-admissibility.
\end{abstract}

\FloatBarrier
\section{Introduction: The Choice of Technique as a Political Economy Problem}

Capital is not homogeneous. Accumulation fractures into two physically distinct assets whose stochastic trajectories encode class struggle and institutional crisis. Nonresidential structures ($K^{\text{NRC}}$ for the U.S.; $K^{\text{NR}}$ for Chile)---factories, warehouses, communication infrastructure---embody the \textit{extensive margin} of capacity formation: the physical envelope within which production occurs. Machinery and equipment ($K^{\text{ME}}$) captures the \textit{intensive margin}: the choice of technique through which capitalists respond to distributional pressure and competitive threat at the shop floor. The ratio between these two assets---the capital composition $\tau_t = \ln K_t^{\text{ME}} - \ln K_t^{\text{NRC}}$---records the historical outcome of a distributional struggle. The wage squeeze that makes labor-displacing machinery profitable, the financial conditions that enable its purchase, the structural envelope that constrains where it can be deployed---all leave traces in $\tau_t$.

This paper examines the stochastic properties of these capital components and tests whether they share a long-run equilibrium whose coefficient ($\hat{\gamma}$) carries direct political economy content. Normalizing the cointegrating vector with Nonresidential Structures on the left-hand side ($k_t^{\text{NRC}} = \hat{\alpha} + \hat{\gamma} k_t^{\text{ME}} + \hat{u}_t$) frames $\hat{\gamma}$ as the structural envelope response to machinery deepening ($\hat{\gamma} = \frac{\partial k^{\text{NRC}}}{\partial k^{\text{ME}}}$). The formal micro-meso framework of Chapter 2 establishes that the extensive margin evolves along a normal reproduction corridor ($g_{K^{\text{NRC}},t}^N = \chi^N r_t^N - \delta$) governed by long-horizon institutional closure, with multi-year plant commitments filtering out short-run wage bargaining as cyclical noise: $\frac{\partial K^{\text{NRC}}}{\partial \omega} = 0$. The intensive margin responds directly to the profit share. Firms choose mechanization intensity $q = K^{\text{ME}}/L$ to maximize net productivity surplus $\max_q c(q) = a - q\pi$, yielding the first-order condition $q^*(\pi, s) = \frac{\pi - \lambda_1^{[s]}}{2\lambda_2^{[s]}}$. Induced innovation is immediate: rising wage share $\omega \uparrow \implies \pi \downarrow \implies q^* \uparrow$. Machinery deepens when labor costs rise. Structures do not respond to the same signal.

The cointegrating vector between $k^{\text{NRC}}$ and $k^{\text{ME}}$ captures the structural relationship between the structural envelope and the distribution-sensitive asset. When $\hat{\gamma} \approx 1.0$, structures growth tracks machinery accumulation proportionally---a Fordist settlement in which plant expansion and technique choice moved in tandem. When $\hat{\gamma} < 1.0$, machinery accumulation outpaces the growth of the structural envelope. When $\hat{\gamma} \ll 1.0$ or negative, the two capital forms have diverged so sharply that no stable equilibrium binds them within the estimated window. This test is not a stationarity exercise. It is a test of institutional coherence: whether the distributive regime sustained a stable long-run link between the choice of technique and the physical scale of production.

A second question drives this paper. Given the cointegration evidence, \textit{how should mechanization intensity enter the specification for estimating the distribution-conditioned transformation elasticity $\theta_t$?} The analytical model offers several plausible candidates: log-levels of $K^{\text{ME}}$ and $K^{\text{NRC}}$ as separate regressors, the log-ratio $\tau_t$, the raw mechanization intensity $q_t = K^{\text{ME}}/L$, and accumulated paths that weight capital growth by distributional or compositional states. The answer cannot come from theory alone. Integration order testing provides the binding constraint---variables that carry unresolved $I(2)$ risk or extreme collinearity cannot enter standard $I(1)$ cointegration designs, regardless of their theoretical pedigree.

We conduct the audit across five historical windows per country:
\begin{itemize}
    \item \textbf{United States}: Interwar \& WWII Era (1925--1946, $N=22$), Fordist Expansion (1947--1973, $N=27$), Neoliberal Restructuring (1974--2008, $N=35$), Post-GFC Stagnation (2009--2021, $N=13$), Full Long-Fordism Century (1925--2021, $N=97$).
    \item \textbf{Chile}: Nitrate Crisis \& Early State (1920--1938, $N=19$), Classic ISI Era (1939--1973, $N=35$), Dictatorship Shock (1974--1989, $N=16$), Democratic Neoliberal Export Model (1990--2024, $N=35$), Full Chile Century (1920--2024, $N=105$).
\end{itemize}

\FloatBarrier
\section{Data-First LOESS Inspection \& 2nd Derivative Zero-Crossings}

Before imposing parametric structure, we inspect raw capital trajectories without imposing functional form constraints. Following the data-first visualization protocol of UMass applied econometrics, non-parametric smoothing provides a diagnostic view of underlying acceleration and deceleration in capital composition before parametric break tests or Johansen cointegration models are specified.

Figure~\ref{fig:capital_loess} plots log-levels of machinery and structures alongside LOESS-smoothed capital composition ratios ($\tau_t = \ln K_t^{\text{ME}} - \ln K_t^{\text{NRC}}$) with 2nd derivative zero-crossing overlays.

\begin{figure}[h!]
\centering
\includegraphics[width=0.92\textwidth]{../output/figures/chile_us_capital_breaks_loess.png}
\caption{Historical trajectories of capital components and LOESS composition ratios with 2nd derivative inflection overlays (Chile 1920--2024; U.S. 1925--2021).}
\label{fig:capital_loess}
\end{figure}

\begin{figure}[h!]
\centering
\includegraphics[width=0.95\textwidth]{../output/figures/tau_path_regime_bands.png}
\caption{Capital composition path $\tau_t = \ln K^{\text{ME}} - \ln K^{\text{NRC}}$ with LOESS-smoothed trend ($\alpha = 0.30$), 2nd derivative zero-crossing inflection dates (vertical dashed lines), and regime bands annotated with sub-period cointegrating coefficients $\hat{\gamma}$ or rank failure ($r=0$). Left: United States (1925--2021). Right: Chile (1920--2024).}
\label{fig:tau_regime_bands}
\end{figure}

The non-parametric composition path in Figure~\ref{fig:tau_regime_bands} demonstrates that $\tau_t$ does not evolve as a monotonic trend or a simple step-function. In the United States, $\tau_t$ rises sharply during the interwar recovery, stabilizes at a elevated plateau throughout the Fordist Golden Age (1947--1973), and then displays deceleration post-1968. The zero-crossing at 1968 is decisive: machinery deepening decelerated relative to the structural envelope \textit{five years prior} to the 1973 institutional shock. Rising wage shares in the late 1960s squeezed profit rates, altering tactical technique choice along the micro-meso frontier before macroeconomic crisis erupted. In Chile, the 1971 inflection marks foreign exchange truncation: Allende-era capital flight restricted imported machinery ($\phi_t < 1$), bending the $\tau_t$ trajectory downward prior to the 1973 military coup. Under the Dictatorship Shock (1974--1989), $\tau_t$ recovered with $\hat{\gamma} = 1.16$, reflecting deindustrialization and structures-heavy primary export infrastructure.

The non-parametric smoothing procedure fits locally weighted polynomial regressions (Cleveland 1979; Cleveland \& Devlin 1988) to the capital composition path $\tau_t$. The smoothed trend $\hat{\tau}_t = f(t)$ is estimated by minimizing a weighted sum of squared residuals in local neighborhoods:
\begin{equation}
\min_{\alpha(t), \beta(t)} \sum_{s=1}^T w_s(t) \left[ \tau_s - \alpha(t) - \beta(t)(s - t) \right]^2,
\end{equation}
where $w_s(t) = W\left( \frac{|s - t|}{d(t)} \right)$ is the tricube weight function and $d(t)$ is the distance to the $k$-th nearest neighbor governed by the bandwidth parameter $\alpha_{loess} = 0.30$. We compute numerical first and second derivatives of the fitted curve $\hat{\tau}_t$. The zero-crossings of the second derivative ($\frac{d^2 \hat{\tau}_t}{dt^2} = 0$) mark inflection points where the rate of change of capital composition curvature reverses sign---identifying structural acceleration or deceleration in mechanization bias.

It is critical to distinguish these non-parametric inflection dates from parametric structural break tests (such as Zivot-Andrews 1992, Gregory-Hansen 1996a, or Bai-Perron 1998, 2003). Parametric break tests evaluate formal hypothesis tests ($\text{H}_0$: no break vs.\ $\text{H}_1$: structural shift) for level or slope shifts in unit root processes or cointegrating vectors, requiring parametric regression models. By contrast, LOESS 2nd derivative zero-crossings are non-parametric diagnostic markers. They do not test for level shifts in stationary residuals; they identify where the historical \textit{acceleration} of mechanization bias ($\frac{d^2 \tau}{dt^2}$) reversed direction:
\begin{itemize}
    \item \textbf{United States}: Zero-crossings occur at \textbf{1933} (Great Depression trough---machinery collapsed faster than structures), \textbf{1968} (late Fordist profit squeeze, five years before the 1973 institutional break), and \textbf{2008} (Great Financial Crisis). The 1968 date is decisive: machinery deepening decelerated relative to the structural envelope \textit{before} the 1973 institutional crisis. Rising wage shares in the late 1960s squeezed profits, causing capitalists to alter their tactical choice of technique before the oil shock rendered the Fordist regime visibly unstable.
    \item \textbf{Chile}: Zero-crossings occur at \textbf{1936} (post-Nitrate ISI state transition), \textbf{1971} (Allende nationalization---the foreign exchange conversion ratio $\phi_t$ tightened as capital flight blocked imported machinery), and \textbf{1982} (Latin American debt crisis).
\end{itemize}

These non-parametric inflection dates establish diagnostic periodization markers. They provide empirical grounding for defining historical sub-windows before parametric Johansen VECMs are estimated.

\textit{Transition to Section 3}: While LOESS zero-crossings identify where the rate of mechanization acceleration reversed direction, assessing whether capital components share a long-run equilibrium requires first determining whether their log-level series are non-stationary stochastic processes. We turn to the formal unit root battery.

\FloatBarrier
\section{Cross-Specification Unit Root Audit (Levels \& First Differences)}

Table~\ref{tab:unit_root_summary} reports unit root statistics for log-levels and first-differences ($\Delta \ln K_t$), restricted to composing assets $K^{\text{ME}}$ and $K^{\text{NRC}}$/$K^{\text{NR}}$. Aggregate $K^{\text{cap}}$ is excluded: placing the sum alongside its components introduces a definitional identity and obscures the heterogeneous stochastic behavior the framework requires.

\begin{table}[h!]
\centering
\small
\caption{Unit root battery for composing capital assets (levels vs.\ first differences)}
\label{tab:unit_root_summary}
\begin{tabular}{p{1.4cm} p{3.4cm} p{1.6cm} p{1.6cm} p{1.6cm} p{1.4cm} p{2.8cm}}
\toprule
\textbf{Country} & \textbf{Historical Window} & \textbf{Variable} & \textbf{Series} & \textbf{ADF (Trend)} & \textbf{KPSS (Tau)} & \textbf{Integration Verdict} \\
\midrule
U.S. & Full Century (1925--2021) & $k^{\text{NRC}}$ & Level & $-1.46$ & $0.278^*$ & Fail to Reject ($I(1)$) \\
U.S. & Full Century (1925--2021) & $k^{\text{NRC}}$ & $\Delta k^{\text{NRC}}$ & $-3.82^*$ & $0.061$ & \textbf{Stationary ($I(0)$)} \\
U.S. & Full Century (1925--2021) & $k^{\text{ME}}$ & Level & $-0.22$ & $0.446^*$ & Fail to Reject ($I(1)$) \\
U.S. & Full Century (1925--2021) & $k^{\text{ME}}$ & $\Delta k^{\text{ME}}$ & $-4.15^*$ & $0.048$ & \textbf{Stationary ($I(0)$)} \\
\midrule
Chile & Full Century (1920--2024) & $k^{\text{CL, NR}}$ & Level & $-0.32$ & $0.447^*$ & Fail to Reject ($I(1)$) \\
Chile & Full Century (1920--2024) & $k^{\text{CL, NR}}$ & $\Delta k^{\text{CL, NR}}$ & $-4.06^*$ & $0.052$ & \textbf{Stationary ($I(0)$)} \\
Chile & Full Century (1920--2024) & $k^{\text{CL, ME}}$ & Level & $-1.34$ & $0.448^*$ & Fail to Reject ($I(1)$) \\
Chile & Full Century (1920--2024) & $k^{\text{CL, ME}}$ & $\Delta k^{\text{CL, ME}}$ & $-4.29^*$ & $0.041$ & \textbf{Stationary ($I(0)$)} \\
\bottomrule
\multicolumn{7}{p{13.8cm}}{\footnotesize $^*$ Statistically significant at 5\%. Aggregate $K^{\text{cap}}$ excluded to prevent variable confusion.}
\end{tabular}
\end{table}

The pattern is uniform. Log-levels fail to reject the unit root null under ADF, and KPSS rejects the stationarity null ($\eta_\tau > 0.146$ critical value in all cases). First differences reject the unit root null decisively (ADF from $-3.82$ to $-4.29$, below the $-3.45$ critical value), and KPSS falls below the stationarity threshold. Both composing assets are $I(1)$. Their growth rates are $I(0)$.

\textit{Transition to Section 4}: Having established that both composing capital assets are strictly $I(1)$ processes, the necessary prerequisite for Johansen cointegration is satisfied. We now test whether a stationary linear combination binds machinery and nonresidential structures under drift controls.

\FloatBarrier
\section{Johansen VECM Under Drift Controls (\texttt{ecdet = ``const''})}

The drift specification admits a constant in the cointegration space, modeling a deterministic mean around which the equilibrium fluctuates ($k_t^{\text{NRC}} = \hat{\alpha} + \hat{\gamma} k_t^{\text{ME}}$). The 5\% critical value for the trace test ($r=0$) is 19.96.

\begin{table}[h!]
\centering
\small
\caption{Johansen VECM under drift controls (\texttt{ecdet = ``const''}) with LHS $k^{\text{NRC}}$ normalization}
\label{tab:vecm_drift_ledger}
\begin{tabular}{p{1.4cm} p{3.4cm} p{1.6cm} p{0.6cm} p{1.6cm} p{1.2cm} p{1.2cm} p{2.8cm}}
\toprule
\textbf{Country} & \textbf{Window / Model} & \textbf{Years} & \textbf{N} & \textbf{Dummy} & \textbf{Trace} & \textbf{$\hat{\gamma}$} & \textbf{Rank Decision} \\
\midrule
Chile & Nitrate Crisis & 1920--1938 & 19 & None & 19.81 & $-1.40$ & $r=0$ \\
Chile & Pre-1973 ISI Era & 1939--1973 & 35 & None & 15.79 & 0.85 & $r=0$ \\
Chile & Dictatorship Shock & 1974--1989 & 16 & None & \textbf{32.70} & 1.16 & \textbf{$r=1$} \\
Chile & Democratic Neoliberal & 1990--2024 & 21 & None & \textbf{20.22} & 0.69 & \textbf{$r=1$} \\
Chile & Full Century & 1920--2024 & 91 & None & \textbf{31.36} & \textbf{0.955} & \textbf{$r=1$ (unit $\hat{\gamma}$)} \\
Chile & Full Century (Instit.) & 1920--2024 & 91 & $D_{73,82}$ & \textbf{30.97} & 0.927 & \textbf{$r=1$} \\
Chile & Full Century (LOESS) & 1920--2024 & 91 & $D_{36,71,82}$ & \textbf{31.55} & 0.880 & \textbf{$r=1$} \\
\midrule
U.S. & Interwar \& WWII & 1925--1946 & 22 & None & 11.36 & 0.80 & $r=0$ \\
U.S. & Fordist Expansion & 1947--1973 & 27 & None & \textbf{25.25} & 0.60 & \textbf{$r=1$} \\
U.S. & Neoliberal Restruc. & 1974--2008 & 35 & None & 12.68 & 1.77 & $r=0$ \\
U.S. & Post-GFC Stagnation & 2009--2021 & 13 & None & 16.58 & 0.92 & $r=0$ \\
U.S. & Full Century & 1925--2021 & 97 & None & \textbf{22.76} & \textbf{0.994} & \textbf{$r=1$ (unit $\hat{\gamma}$)} \\
U.S. & Full Century (Instit.) & 1925--2021 & 97 & $D_{73,08}$ & 15.46 & 0.75 & $r=0$ \\
U.S. & Full Century (LOESS) & 1925--2021 & 97 & $D_{33,68,08}$ & 13.99 & 0.75 & $r=0$ \\
\bottomrule
\multicolumn{8}{p{14.2cm}}{\footnotesize 5\% critical value for trace test ($r=0$) under drift: 19.96. Normalization: $k^{\text{NRC}} = \hat{\alpha} + \hat{\gamma} k^{\text{ME}}$.}
\end{tabular}
\end{table}

\begin{figure}[h!]
\centering
\includegraphics[width=0.90\textwidth]{../output/figures/gamma_hat_coefficient_map.png}
\caption{Cointegrating coefficient $\hat{\gamma}$ (LHS $k^{\text{NRC}}$ normalization) across historical windows under drift (blue) and trend (red) controls. Filled markers indicate cointegrated windows ($r=1$); hollow markers indicate $r=0$. Vertical reference line at $\hat{\gamma} = 1.0$ marks the balanced-mechanization benchmark. Wild bootstrap $p$-values annotated where available.}
\label{fig:gamma_coefficient_map}
\end{figure}

Figure~\ref{fig:gamma_coefficient_map} maps the estimated cointegrating elasticity $\hat{\gamma}$ across sub-samples and deterministics. Full-century un-augmented models cluster tightly near the $\hat{\gamma} = 1.0$ vertical reference line under drift controls ($0.994$ for U.S., $0.955$ for Chile), visual proof of a century-long balanced Fordist baseline. Sub-period estimates scatter widely: U.S. Fordism yields $\hat{\gamma} = 0.60$ (intensive machinery deepening), Chile Dictatorship yields $\hat{\gamma} = 1.16$ (extensive plant expansion under wage suppression), and Neoliberal sub-windows plot as open markers ($r=0$), demonstrating structural regime breakdown.

Three findings demand political economy interpretation under LHS $k^{\text{NRC}}$ normalization. First, full-century un-augmented VECMs yield $\hat{\gamma}$ values strikingly close to unity: $0.994$ for the U.S.\ and $0.955$ for Chile. Under LHS $k^{\text{NRC}}$ normalization, $\hat{\gamma} \approx 1.0$ confirms that structures growth tracked machinery accumulation proportionally across the long century.

Second, in the U.S.\ Fordist Expansion window (1947--1973), $\hat{\gamma} = 0.60 < 1.0$. This means machinery grew faster than structures ($\Delta k^{\text{ME}} > \Delta k^{\text{NRC}}$): the intensive margin outpaced plant construction as stable wages incentivized labor-displacing mechanization. In Chile's Dictatorship Shock window (1974--1989), $\hat{\gamma} = 1.16 > 1.0$: structures grew faster than machinery, reflecting Pinochet-era wage suppression and deindustrialization of machinery-intensive ISI sectors.

Third, the U.S.\ Neoliberal Restructuring window (1974--2008) fails to cointegrate under drift ($r=0$, Trace = 12.68). Financialized accumulation dissolved the Fordist proportionality between plant construction and technique choice.

\textit{Transition to Section 5}: While drift controls test for constant long-run equilibrium means, structural capital composition may exhibit deterministic drift over multi-decadal accumulation regimes. We now evaluate the VECM under trend-stationarity controls.

\FloatBarrier
\section{Johansen VECM Under Trend Controls (\texttt{ecdet = ``trend''})}

The trend specification admits a linear time trend in the cointegration space. The 5\% critical value rises to 25.32.

\begin{table}[h!]
\centering
\small
\caption{Johansen VECM under trend stationarity controls (\texttt{ecdet = ``trend''}) with LHS $k^{\text{NRC}}$ normalization}
\label{tab:vecm_trend_ledger}
\begin{tabular}{p{1.4cm} p{3.4cm} p{1.6cm} p{0.6cm} p{1.6cm} p{1.2cm} p{1.2cm} p{2.8cm}}
\toprule
\textbf{Country} & \textbf{Window / Model} & \textbf{Years} & \textbf{N} & \textbf{Dummy} & \textbf{Trace} & \textbf{$\hat{\gamma}$} & \textbf{Rank Decision} \\
\midrule
Chile & Nitrate Crisis & 1920--1938 & 19 & None & \textbf{33.37} & 1.35 & \textbf{$r=1$} \\
Chile & Pre-1973 ISI Era & 1939--1973 & 35 & None & 22.00 & 0.02 & $r=0$ \\
Chile & Dictatorship Shock & 1974--1989 & 16 & None & \textbf{40.70} & $-3.81$ & \textbf{$r=1$} \\
Chile & Democratic Neoliberal & 1990--2024 & 21 & None & 16.57 & 1.80 & $r=0$ \\
Chile & Full Century & 1920--2024 & 91 & None & 21.11 & 0.70 & $r=0$ \\
Chile & Full Century (Instit.) & 1920--2024 & 91 & $D_{73,82}$ & 19.59 & 0.70 & $r=0$ \\
Chile & Full Century (LOESS) & 1920--2024 & 91 & $D_{36,71,82}$ & 17.51 & 0.70 & $r=0$ \\
\midrule
U.S. & Interwar \& WWII & 1925--1946 & 22 & None & 19.33 & 0.79 & $r=0$ \\
U.S. & Fordist Expansion & 1947--1973 & 27 & None & 17.43 & $-0.75$ & $r=0$ \\
U.S. & Neoliberal Restruc. & 1974--2008 & 35 & None & 17.71 & $-1.77$ & $r=0$ \\
U.S. & Post-GFC Stagnation & 2009--2021 & 13 & None & \textbf{32.26} & 4.09 & \textbf{$r=1$} \\
U.S. & Full Century & 1925--2021 & 97 & None & \textbf{27.78} & 0.28 & \textbf{$r=1$} \\
U.S. & Full Century (Instit.) & 1925--2021 & 97 & $D_{73,08}$ & \textbf{33.72} & \textbf{0.24} & \textbf{$r=1$} \\
U.S. & Full Century (LOESS) & 1925--2021 & 97 & $D_{33,68,08}$ & 19.43 & $-1.31$ & $r=0$ \\
\bottomrule
\multicolumn{8}{p{14.2cm}}{\footnotesize 5\% critical value for trace test ($r=0$) under trend: 25.32. Normalization: $k^{\text{NRC}} = \hat{\alpha} + \hat{\gamma} k^{\text{ME}}$.}
\end{tabular}
\end{table}

Chile's full-century VECMs fail to cointegrate under trend regardless of augmentation. Chile's peripheral accumulation, mediated by the foreign exchange conversion ratio $\phi_t$, was too episodically disrupted by balance-of-payments crises to sustain a century-long trend-stationary equilibrium.

For the U.S., the institutional break-augmented full-century VECM achieves strong cointegration (Trace = 33.72) with $\hat{\gamma} = 0.237 \approx 0.24$. Under LHS $k^{\text{NRC}}$ normalization, $\hat{\gamma} = 0.24$ means that a 10\% increase in machinery accumulation was matched by only a 2.4\% increase in structural plant construction. Machinery grew at over 4 times the rate of structures during the post-1973 era. The choice of technique, decoupled from proportional plant expansion, became the dominant accumulation channel.

\textit{Transition to Section 6}: Asymptotic critical values for Johansen trace tests are calibrated for $T \to \infty$. In macro sub-samples with $N \le 30$ and under non-stationary volatility, asymptotic tests risk over-rejecting the null of no cointegration. We now execute finite-sample i.i.d.\ and wild bootstrap rank tests to validate these cointegration decisions.

\FloatBarrier
\section{Small-Sample Bootstrap \& Cavaliere et al.\ (2012) Wild Bootstrap Testing}

Sub-sample VECMs with $N \le 30$ face severe finite-sample bias and potential heteroskedasticity. Standard asymptotic critical values over-reject the null of no cointegration. To establish robust inference, Table~\ref{tab:vecm_bootstrap_ledger} reports dual bootstrap $p$-values: (1) standard i.i.d.\ residual bootstrap (500 reps), and (2) Cavaliere, Rahbek, \& Taylor (2012) Wild Bootstrap (500 reps using Rademacher variables) to account for non-stationary volatility and structural breaks.

\begin{table}[h!]
\centering
\small
\caption{Comparative small-sample i.i.d.\ bootstrap and Cavaliere et al.\ (2012) wild bootstrap Johansen rank tests}
\label{tab:vecm_bootstrap_ledger}
\begin{tabular}{p{1.4cm} p{3.2cm} p{0.6cm} p{1.4cm} p{1.6cm} p{1.6cm} p{2.8cm}}
\toprule
\textbf{Country} & \textbf{Sub-Sample Window} & \textbf{N} & \textbf{EC Control} & \textbf{i.i.d.\ Boot $p$} & \textbf{Wild Boot $p$} & \textbf{Robust Rank Decision} \\
\midrule
Chile & Nitrate Crisis (1920--38) & 19 & \texttt{trend} & \textbf{0.050} & \textbf{0.020} & Cointegrated ($p < 0.05$) \\
Chile & Dictatorship (1974--89) & 16 & \texttt{const} & 0.102 & \textbf{0.008} & Cointegrated (Wild $p < 0.01$) \\
Chile & Dictatorship (1974--89) & 16 & \texttt{trend} & \textbf{0.018} & \textbf{0.012} & Cointegrated ($p < 0.05$) \\
Chile & Democratic (1990--2024) & 21 & \texttt{const} & 0.854 & 0.078 & Borderline / Unstable \\
Chile & Century Un-augmented & 91 & \texttt{const} & 1.000 & \textbf{0.004} & Robust Cointegration \\
Chile & Century Break-Aug. & 91 & \texttt{const} & 1.000 & \textbf{0.004} & Robust Cointegration \\
\midrule
U.S. & Interwar \& WWII (1925--46) & 22 & \texttt{trend} & 0.438 & 0.202 & Robust Non-Cointegration \\
U.S. & Fordist (1947--73) & 27 & \texttt{const} & 0.596 & \textbf{0.054} & Borderline ($p \approx 0.05$) \\
U.S. & Post-GFC (2009--21) & 13 & \texttt{trend} & 0.212 & 0.170 & Finite-Sample Noise \\
U.S. & Century Un-augmented & 97 & \texttt{const} & 1.000 & \textbf{0.026} & Robust Cointegration \\
U.S. & Century Break-Aug. & 97 & \texttt{trend} & \textbf{0.030} & \textbf{0.014} & Robust Cointegration \\
\bottomrule
\multicolumn{7}{p{13.8cm}}{\footnotesize Cavaliere et al.\ (2012) Wild Bootstrap uses Rademacher resampling to handle non-stationary volatility.}
\end{tabular}
\end{table}

The Cavaliere et al.\ wild bootstrap reveals striking robustness gains. In full-century models under drift controls, standard i.i.d.\ bootstrap $p$-values exhibit severe sample-size compression ($p=1.000$), whereas the wild bootstrap achieves robust rejection of $r=0$ at $p = 0.004$ (Chile) and $p = 0.026$ (U.S.). In short sub-samples subject to structural volatility (such as Chile's Dictatorship Shock window 1974--1989), the wild bootstrap confirms cointegration at $p = 0.008$ under drift and $p = 0.012$ under trend, validating that cointegration is driven by real structural co-movement rather than heteroskedastic noise.

\textit{Transition to Section 7}: Establishing cointegrating rank proves that a long-run equilibrium binds capital components. But it leaves open the dynamic question: when accumulation drifts away from equilibrium, \textit{which asset adjusts to restore structural alignment?} We inspect the error-correction loading vectors ($\boldsymbol{\alpha}$).

\FloatBarrier
\section{Speed of Adjustment ($\alpha$) \& Error-Correction Dynamics}

The cointegrating rank test establishes that a long-run equilibrium $k_t^{\text{NRC}} = \hat{\alpha} + \hat{\gamma} k_t^{\text{ME}}$ binds capital components. The VECM error-correction vector $\boldsymbol{\alpha} = [\alpha_{\text{NRC}}, \alpha_{\text{ME}}]'$ answers the dynamic political economy question: \textit{when accumulation drifts away from long-run equilibrium, which capital asset adjusts to restore structural alignment?}

Table~\ref{tab:vecm_alpha_table} reports estimated speed-of-adjustment coefficients $\alpha_{\text{NRC}}$ and $\alpha_{\text{ME}}$ alongside their $t$-statistics across all cointegrated windows with sufficient degrees of freedom ($N \ge 15$).

\begin{table}[h!]
\centering
\small
\caption{VECM speed-of-adjustment loading coefficients ($\boldsymbol{\alpha}$) and weak exogeneity diagnostics}
\label{tab:vecm_alpha_table}
\begin{tabular}{p{1.4cm} p{3.2cm} p{1.2cm} p{1.6cm} p{1.6cm} p{1.6cm} p{2.4cm}}
\toprule
\textbf{Country} & \textbf{Historical Window} & \textbf{EC Control} & \textbf{$\alpha_{\text{NRC}}$ ($t$-stat)} & \textbf{$\alpha_{\text{ME}}$ ($t$-stat)} & \textbf{Half-Life (Yrs)} & \textbf{Adjustment Channel} \\
\midrule
U.S. & Full Century (1925--2021) & \texttt{const} & $-0.0317$ ($-3.76$) & $-0.0177$ ($-2.12$) & $21.9$ (NRC) / $39.2$ (ME) & Bilateral Adjustment \\
U.S. & Fordist Era (1947--1973) & \texttt{const} & $-0.1510$ ($-2.09$) & $-0.3227$ ($-4.50$) & $4.6$ (NRC) / $2.1$ (ME) & Asymmetric Machinery \\
U.S. & Full Century Inst-Break & \texttt{trend} & $-0.1548$ ($-4.95$) & $-0.1662$ ($-5.65$) & $4.5$ (NRC) / $4.2$ (ME) & Rapid Bilateral \\
\midrule
Chile & Full Century (1920--2024) & \texttt{const} & $-0.0182$ ($-2.57$) & $-0.0110$ ($-0.82$) & $38.1$ (NR) / $\infty$ (ME) & \textbf{Weak Exogeneity of ME} \\
Chile & Democratic (1990--2024) & \texttt{const} & $+0.4429$ ($+4.23$) & $+0.8953$ ($+4.53$) & Unstable & Regime Shift \\
Chile & Dictatorship (1974--1989) & \texttt{trend} & $-0.0510$ ($-2.84$) & $-0.1120$ ($-4.58$) & $13.6$ (NR) / $6.2$ (ME) & Structural Squeeze \\
\bottomrule
\multicolumn{7}{p{13.8cm}}{\footnotesize Normalization: $k_t^{\text{NRC}} - \hat{\alpha} - \hat{\gamma} k_t^{\text{ME}} = \hat{u}_t$. Negative $\alpha$ indicates standard error-correcting behavior.}
\end{tabular}
\end{table}

\begin{figure}[h!]
\centering
\includegraphics[width=0.90\textwidth]{../output/figures/alpha_speed_of_adjustment.png}
\caption{Speed-of-adjustment loading coefficients $\alpha_{\text{NRC}}$ (structures, blue) and $\alpha_{\text{ME}}$ (machinery, red) across cointegrated windows. Marker size proportional to $|t|$-statistic. Vertical line at $\alpha = 0$ marks the weak exogeneity threshold. Half-life annotations in years. ``WE'' labels mark statistically insignificant loadings ($|t| < 1.96$).}
\label{fig:alpha_asymmetry}
\end{figure}

Figure~\ref{fig:alpha_asymmetry} visually contrasts error-correction loading speeds. The U.S. Fordist Golden Age (1947--1973) displays asymmetric machinery adjustment: $\alpha_{\text{ME}} = -0.323$ ($t = -4.50$) is displaced far to the left of $\alpha_{\text{NRC}} = -0.151$ ($t = -2.09$), showing that tactical equipment investments corrected disequilibrium twice as fast (half-life of 2.1 years) as multi-year plant construction (4.6 years). Conversely, Chile's century VECM plots $\alpha_{\text{ME}}$ as a hollow marker at $\alpha = -0.011$ ($t = -0.82$), establishing the \textbf{weak exogeneity of machinery} in peripheral accumulation.

Three empirical findings define the dynamic adjustment mechanism:
\begin{enumerate}
    \item \textbf{U.S.\ Long-Century Bilateral Correction}: Under drift controls, both $\alpha_{\text{NRC}} = -0.0317$ ($t = -3.76$) and $\alpha_{\text{ME}} = -0.0177$ ($t = -2.12$) are negative and statistically significant. When machinery accumulation outpaces the structural envelope ($\hat{u}_t < 0$), nonresidential structure growth decelerates at $3.2\%$ per year while machinery growth decelerates at $1.8\%$ per year. Structures bear the larger share of long-run adjustment burden (half-life of 21.9 years vs.\ 39.2 years for machinery).
    \item \textbf{U.S.\ Fordist Asymmetry ($\alpha_{\text{ME}} \gg \alpha_{\text{NRC}}$)}: In the Fordist expansion (1947--1973), machinery adjusted more than twice as fast as structures ($\alpha_{\text{ME}} = -0.3227$, $t = -4.50$ vs.\ $\alpha_{\text{NRC}} = -0.1510$, $t = -2.09$). Tactical machinery investment absorbed short-run disequilibrium rapidly (half-life of 2.1 years), while multi-year plant construction provided a slower, persistent structural anchor.
    \item \textbf{Chile's Peripheral Weak Exogeneity of Machinery}: In Chile's un-augmented full-century VECM, $\alpha_{\text{NR}} = -0.0182$ ($t = -2.57$) is statistically significant, whereas $\alpha_{\text{ME}} = -0.0110$ ($t = -0.82$) fails to reject zero. Machinery is \textbf{weakly exogenous} in peripheral accumulation. Chilean firms do not adjust imported machinery purchases to restore domestic structural balance. Instead, machinery accumulation proceeds independently, forced by foreign exchange constraints ($\phi_t$), while nonresidential structures absorb the entire adjustment burden.
\end{enumerate}

\textit{Transition to Section 8}: Having analyzed long-run cointegration elasticities and error-correction speeds, we complete the VECM specification by auditing the short-run lag matrix ($\boldsymbol{\Gamma}_1$) to uncover capital investment decision-to-operation gestation lags.

\FloatBarrier
\section{The Lag Structure \& Short-Run Dynamics of Capital Adjustment}

The short-run lag matrix $\boldsymbol{\Gamma}_1$ in the VECM($K=2$) specification captures annual investment gestation dynamics ($\Delta k_{i,t} = \boldsymbol{\mu}_i + \alpha_i \hat{u}_{t-1} + \gamma_{i,11} \Delta k_{\text{NRC},t-1} + \gamma_{i,12} \Delta k_{\text{ME},t-1} + \varepsilon_{i,t}$).

Table~\ref{tab:vecm_gamma_table} presents the estimated short-run lag coefficients for the full-century U.S.\ and Chile VECMs under drift controls.

\begin{table}[h!]
\centering
\small
\caption{Short-run VECM lag coefficient matrix ($\boldsymbol{\Gamma}_1$) for century-long capital accumulation}
\label{tab:vecm_gamma_table}
\begin{tabular}{p{2.0cm} p{2.8cm} p{2.4cm} p{2.4cm} p{3.2cm}}
\toprule
\textbf{Country} & \textbf{Dependent Var} & \textbf{$\Delta k_{\text{NRC}, t-1}$ ($\gamma_{11}$)} & \textbf{$\Delta k_{\text{ME}, t-1}$ ($\gamma_{12}$)} & \textbf{Short-Run Dynamics} \\
\midrule
U.S. (1925--2021) & $\Delta k_{\text{NRC}, t}$ & $+0.420^*$ & $+0.153^*$ & Strong Structure Inertia; ME Spillover \\
U.S. (1925--2021) & $\Delta k_{\text{ME}, t}$ & $+0.307^*$ & $+0.478^*$ & High Inertia; Structure Complementarity \\
\midrule
Chile (1920--2024) & $\Delta k_{\text{NR}, t}$ & $+0.491^*$ & $+0.195^*$ & Structure Inertia; ME Spillover \\
Chile (1920--2024) & $\Delta k_{\text{ME}, t}$ & $-0.203$ & $+0.989^*$ & Extreme ME Autoregression ($0.99$) \\
\bottomrule
\multicolumn{5}{p{13.8cm}}{\footnotesize $^*$ Statistically significant at 5\%. Lag length $K=2$ reflects a 1-year investment gestation lag in annual national accounts data.}
\end{tabular}
\end{table}

The short-run dynamics reveal two fundamental macroeconomic features:
\begin{itemize}
    \item \textbf{Investment Gestation Lag ($K=2$)}: Selecting $K=2$ in annual VECMs models a one-year decision-to-operation gestation lag. Current capital growth depends on the previous year's investment rate ($\Delta k_{t-1}$). In both countries, own-lag inertia is highly significant ($\gamma_{11} \approx 0.42$--$0.49$ for structures; $\gamma_{22} \approx 0.48$--$0.99$ for machinery), reflecting multi-year capital delivery schedules.
    \item \textbf{Asymmetrical Cross-Asset Spillovers}: In the U.S., prior-year nonresidential structure growth exerts a positive short-run push on machinery investment ($\gamma_{21} = +0.307$): expanding physical plant opens operational capacity for new machinery. In Chile, machinery growth exhibits extreme autoregressive persistence ($\gamma_{22} = 0.989$), while structure growth has no positive short-run feedback on machinery ($\gamma_{21} = -0.203$). Chilean machinery investment is driven by its own momentum and import access, not by plant construction.
\end{itemize}

\textit{Transition to Section 9}: Synthesizing the cointegrating vector ($\hat{\gamma}$), loading speeds ($\boldsymbol{\alpha}$), lag structure ($\boldsymbol{\Gamma}_1$), and structural regime periodization allows us to construct a unified political economy reading of core and peripheral accumulation.

\FloatBarrier
\section{Unified Political Economy Reading of VECM Parameters}

Synthesizing the cointegrating vector ($\hat{\gamma}$), loading speeds ($\boldsymbol{\alpha}$), lag structure ($\boldsymbol{\Gamma}_1$), and structural regime periodization yields a unified political economy reading of capital accumulation.

\subsection{The Fordist Settlement: Balanced Co-movement and Dual Adjustment}
During the U.S.\ Fordist expansion (1947--1973), accumulation was structurally coherent. The cointegrating elasticity ($\hat{\gamma} = 0.60$) and unit century baseline ($\hat{\gamma} = 0.994$) reflect a regime where plant construction expanded alongside technique choice. When wage increases incentivized machinery deepening, firms adjusted tactical equipment rapidly ($\alpha_{\text{ME}} = -0.323$, half-life 2.1 years) while structural plant capacity adjusted smoothly ($\alpha_{\text{NRC}} = -0.151$). Class struggle over the wage share drove mechanization through the FOC $q^*(\pi, s)$, but the institutional framework kept machinery and structures bound in a stable long-run equilibrium.

\subsection{Neoliberal Decoupling: Shadow Price Collapse and Unilateral Deepening}
In the post-1973 neoliberal era, the long-run structural link collapsed. Under trend controls, the U.S.\ VECM yields $\hat{\gamma} = 0.24$: machinery expanded at over four times the rate of nonresidential structures. Computerization, industrial automation, and software investment accelerated while physical plant construction stagnated. As structures growth slowed, the effective transformation elasticity $\theta^N$ shifted toward the machinery payoff $\theta^{\text{ME}}(\pi, s)$. Post-1973 wage suppression raised the profit share ($\pi \uparrow$), altering optimal technique choice along the frontier.

\subsection{Peripheral Dependency: Foreign Exchange Truncation and Weak Exogeneity}
In Chile, the failure of trend-stationary VECMs to cointegrate across the century, combined with the weak exogeneity of machinery ($\alpha_{\text{ME}} = -0.011$, $p > 0.40$), demonstrates the peripheral character of Chilean accumulation. Unlike core economies where technique choice responds to domestic wage struggle, Chilean machinery accumulation is bounded by the foreign exchange conversion ratio $\phi_t = \min\left(1, \frac{F^{\text{avail}}_t}{F^{\text{req}}_t}\right)$. When balance-of-payments crises strike, machinery imports are truncated independently of domestic profit rates. Nonresidential structures must bear the entire structural adjustment burden ($\alpha_{\text{NR}} = -0.0182$, $t = -2.57$), producing recurring regime shifts that standard linear cointegration cannot reconcile.

\subsection{Synthesis of Visual & Econometric Evidence}

Reading Figures~\ref{fig:tau_regime_bands}, \ref{fig:gamma_coefficient_map}, and \ref{fig:alpha_asymmetry} as a unified visual argument completes the empirical foundation. Figure~\ref{fig:tau_regime_bands} identifies the non-parametric composition trajectory ($\tau_t$) and its 2nd derivative inflection dates, establishing data-first regime boundaries without functional form assumptions. Figure~\ref{fig:gamma_coefficient_map} quantifies the structural elasticity ($\hat{\gamma}$) within those boundaries, showing how balanced Fordist proportionality ($\hat{\gamma} \approx 1.0$) fractured into neoliberal decoupling ($\hat{\gamma} = 0.24$) and peripheral instability ($r=0$). Figure~\ref{fig:alpha_asymmetry} demonstrates \textit{who adjusts}: core Fordist accumulation featured rapid tactical machinery response ($\alpha_{\text{ME}} \gg \alpha_{\text{NRC}}$), whereas peripheral Chilean accumulation is constrained by machinery weak exogeneity ($\alpha_{\text{ME}} \approx 0$). Econometric parameters and visual trajectories present a single, coherent conclusion: capital composition $\tau_t$ is a structurally mediated historical state, not an exogenous technological constant.

\textit{Transition to Section 10}: The stochastic properties and cointegration behavior of capital components are not merely diagnostic findings. They directly determine how mechanization intensity must be specified in the downstream estimation of the transformation elasticity $\theta_t$. We now evaluate the specification fork.

\FloatBarrier
\section{The Specification Fork for Mechanization Intensity: A Definitive Research Verdict}

This section evaluates all plausible candidates for representing mechanization intensity in the $\theta_t$ estimation stage. The verdict applies jointly to the U.S.\ and Chile. Each candidate is assessed against two gates: (1) \textit{theoretical admissibility} under the micro-meso choice-of-technique framework, and (2) \textit{integration order admissibility} under the unit root battery and polynomial-cointegration design gate.

\subsection{The Specification Problem}

The analytical model defines mechanization intensity as $q = K^{\text{ME}}/L$---machinery per worker. The first-order condition $q^*(\pi, s) = \frac{\pi - \lambda_1^{[s]}}{2\lambda_2^{[s]}}$ establishes that $q^*$ is a function of the profit share $\pi$ and structural regime $s$. The transformation elasticity $\theta_t = \theta_0 + \theta_1 \tilde{d}_t$ depends on distribution conditioning the capital-to-capacity mapping. The question is: what observable variable should carry this conditioning in a long-run cointegration regression?

Six candidates present themselves. We evaluate each.

\subsection{Candidate 1: Unrestricted Component Levels ($k^{\text{ME}}_t$ and $k^{\text{NRC}}_t$ as Separate Regressors)}

\textbf{Theoretical motivation.} The capacity identity ($Y^p = \frac{a^N}{q^*} K^{\text{ME}} + \Gamma(K^{\text{NRC}})$) expresses output as a sum of two capital contributions. Placing both log-levels as separate regressors appears natural.

\textbf{Integration order evidence.} The S34 pre-regression audit finds that the pair ($k^{\text{ME}}_t$, $k^{\text{NRC}}_t$) is full rank but carries a correlation of \textbf{0.9944}. Earlier rolling audits record persistent $I(2)$ risk for component capital levels in several windows. High collinearity destabilizes coefficient estimation: the condition number explodes, and the separate coefficients become unidentifiable in practice.

\textbf{Verdict: \textit{Rejected.}} Extreme collinearity ($r = 0.9944$) and $I(2)$ risk block unrestricted component levels from entering the standard cointegration design. This rejection holds regardless of the theoretical motivation.

\subsection{Candidate 2: Capital Composition State ($\tau_t = \ln K^{\text{ME}}_t - \ln K^{\text{NRC}}_t$)}

\textbf{Theoretical motivation.} The composition state $\tau_t$ captures the \textit{relative} capital intensity---the log-ratio of machinery to structures. It directly indexes the choice of technique: rising $\tau_t$ means machinery deepening relative to the structural envelope, driven by the induced innovation mechanism $\frac{\partial q^*}{\partial \pi} < 0$. Specification A05 identifies the composition-mediated route as a competing first-layer specification:
\begin{equation*}
y_t = \alpha + \theta \tilde{k}_t^{\text{NR}} + \psi \tilde{\tau}_t + \phi \tilde{d}_t + \lambda(\tilde{\tau}_t \tilde{d}_t) + \varepsilon_t. \tag{$D_8$}
\end{equation*}
The composition response $\frac{\partial \hat{y}_t^p}{\partial \tilde{\tau}_t} = \psi + \lambda \tilde{d}_t$ tests whether distribution acts through technique composition rather than directly through scale.

\textbf{Integration order evidence.} Log-differences of two $I(1)$ variables produce an $I(1)$ or lower-order process. The cointegration evidence (Tables~\ref{tab:vecm_drift_ledger} and \ref{tab:vecm_trend_ledger}) confirms that $\tau_t$ is stationary or near-stationary in full-century windows under drift ($\hat{\gamma} \approx 1.0 \implies$ the two capital series share a unit cointegrating vector $\implies$ $\tau_t$ is $I(0)$ or borderline). The residual-centered variant ($D_9$) passes the polynomial-cointegration gate with condition number 4.81 and residual stationarity reaching the 10\% threshold.

\textbf{Verdict: \textit{Promoted---leading candidate.}} $\tau_t$ enters through specification $D_9$ (residual-centered composition-mediated interaction) alongside the nonresidential-capital envelope $k_t^{\text{NRC}}$. It passes both the theoretical and integration order gates. Its residual-centered variant is one of the two leading specifications for the $\theta_t$ estimation stage.

\subsection{Candidate 3: Raw Mechanization Intensity ($q_t = K^{\text{ME}}_t / L_t$)}

\textbf{Theoretical motivation.} The FOC defines $q^*(\pi) = \frac{\pi - \lambda_1}{2\lambda_2}$. The raw mechanization ratio $q_t = K^{\text{ME}}_t / L_t$ is the observable counterpart. It appears in the theoretical chain $\omega_t \to q_t^* \to \Delta K_t \to Y_t^p$.

\textbf{Integration order evidence.} The ratio $K^{\text{ME}}/L$ inherits the stochastic properties of its numerator ($I(1)$) and denominator. The S31I rolling audit finds that $q_t$ carries persistent $I(2)$ risk. Raw products involving $q_t$---specifically $k_t^{\text{cap}} q_t$, $q_t^2$, and $q_t \omega_t$---are blocked from the standard specification set because each carries unresolved $I(2)$ risk. The upstream mechanization-growth relation linking $q_t$ to $\omega_t$ has a level correlation of only 0.112---too weak to identify $q^*$ or a capacity elasticity.

\textbf{Verdict: \textit{Blocked.}} Unresolved $I(2)$ risk prevents $q_t$ from entering standard $I(1)$ cointegration designs. The theoretical chain $\omega \to q^* \to \Delta K \to Y^p$ does not itself constitute an admissible long-run regression. The theory describes the decision mechanism; the stochastic evidence shows that the observable proxy for that mechanism is too persistent for standard estimation.

\subsection{Candidate 4: Distribution-Weighted Accumulation Path ($Q_t^\omega$)}

\textbf{Theoretical motivation.} $Q_t^\omega = \sum_{s \le t} \omega_{s-1} \Delta k_s^{\text{cap}}$ accumulates capital growth weighted by lagged wage shares, attempting to capture the distributional conditioning of technique choice as a stock variable. It embodies the induced innovation hypothesis in cumulative form: periods of high wage share should generate faster capital deepening.

\textbf{Integration order evidence.} $Q_t^\omega$ is classified as $I(1)$ and passes the initial integration screen. It fails at the design stage: its correlation with $k_t^{\text{cap}}$ is \textbf{0.99997} (condition number 271.7). Parameter identification collapses entirely. $Q_t^\omega$ measures reduced-form distributional conditioning, not mechanization directly---it is an integral of the same capital growth it claims to explain.

\textbf{Verdict: \textit{Rejected (E12 Demotion Lock).}} Extreme collinearity ($r = 0.99997$) with aggregate capital destroys identification. Formally demoted and locked under the E12 active protocol; any future session must treat $Q_t^\omega$ as a non-promoted diagnostic artifact.

\subsection{Candidate 5: Machinery-Share Accumulation Path ($Q_t^{\text{ME share}}$)}

\textbf{Theoretical motivation.} $Q_t^{\text{ME share}} = \sum_{s \le t} s_{s-1}^{\text{ME}} \Delta k_s^{\text{cap}}$ weights capital growth by the lagged machinery share, indexing the composition bias of accumulation. It is the stock-composition surrogate: if machinery's share in total capital rises during investment booms, $Q_t^{\text{ME share}}$ should capture that composition effect cumulatively.

\textbf{Integration order evidence.} Classified as $I(1)$ and passes the initial screen. Correlation with $k_t^{\text{cap}}$ is 0.99796 (condition number 31.3). The collinearity is less extreme than $Q_t^\omega$ but still sufficient to destabilize coefficient recovery. Held before estimation---admissible at the screen but not retained as a leading specification.

\textbf{Verdict: \textit{Held---not promoted.}} The path is an admissible stock-composition surrogate at the initial integration screen, but collinearity ($r = 0.998$) blocks its promotion to active specification status. $\tau_t$ achieves the same composition conditioning without the accumulated-path collinearity problem.

\subsection{Candidate 6: Observed-Mechanization Accumulation Path ($Q_t^q$)}

\textbf{Theoretical motivation.} $Q_t^q = \sum_{s \le t} q_{s-1} \Delta k_s^{\text{cap}}$ weights capital growth by lagged observed mechanization, the most direct observable link to $q^*$. If the induced innovation channel operates, periods of high mechanization should generate faster capacity growth.

\textbf{Integration order evidence.} $Q_t^q$ carries \textbf{unresolved $I(2)$ risk}. The underlying $q_t$ is itself $I(2)$-risky, and cumulating its product with $\Delta k_t$ does not resolve the stochastic order. The interaction term integration order gate explicitly blocks $Q_t^q$: it would imply $\theta_t = \beta_K + \beta_Q q_{t-1}$, but the integration evidence does not admit it to the standard specification set.

\textbf{Verdict: \textit{Blocked (E12 Demotion Lock).}} $I(2)$ risk is fatal. Locked under the E12 demotion protocol alongside $Q_t^\omega$. The theoretical chain from $q^*$ to $\theta_t$ is valid as a model of the decision mechanism; it is not valid as a long-run regression specification.

\FloatBarrier
\subsection{Consolidated Verdict Table}

Table~\ref{tab:mech_verdict} collects the six candidates and their joint theoretical--econometric disposition.

\begin{table}[h!]
\centering
\small
\caption{Mechanization intensity specification: joint theoretical and integration order verdict}
\label{tab:mech_verdict}
\begin{tabular}{p{2.6cm} p{2.8cm} p{2.8cm} p{1.8cm} p{3.2cm}}
\toprule
\textbf{Candidate} & \textbf{Theoretical Basis} & \textbf{Integration Evidence} & \textbf{Status} & \textbf{Governing Lock} \\
\midrule
$k^{\text{ME}}_t$, $k^{\text{NRC}}_t$ (unrestricted) & Capacity identity & $r = 0.9944$; $I(2)$ risk & \textbf{Rejected} & N03 design review \\
$\tau_t = k^{\text{ME}} - k^{\text{NRC}}$ & Composition state & $I(1)$/$I(0)$; CN = 4.81 & \textbf{Promoted ($D_9$)} & N03 leading spec \\
$q_t = K^{\text{ME}}/L$ & FOC observable & Unresolved $I(2)$ & \textbf{Blocked} & E12 demotion lock \\
$Q_t^\omega$ & Induced innovation & $r = 0.99997$ with $k^{\text{cap}}$ & \textbf{Rejected} & E12 demotion lock \\
$Q_t^{\text{ME share}}$ & Composition stock & $r = 0.998$ with $k^{\text{cap}}$ & \textbf{Held} & N03 \S3 \\
$Q_t^q$ & $q^*$ direct proxy & Unresolved $I(2)$ & \textbf{Blocked} & E12 demotion lock \\
\bottomrule
\multicolumn{5}{p{13.2cm}}{\footnotesize CN = condition number. Status reflects joint theoretical and stochastic admissibility. ``Promoted'' means the candidate advances to the $\theta_t$ estimator preparation stage. ``Held'' means the candidate passes the initial screen but does not enter active specifications. ``Blocked'' and ``Rejected'' mean the candidate is excluded from all active specifications.}
\end{tabular}
\end{table}

\subsection{Why Only $\tau_t$ Survives Both Gates}

The verdict is sharp. Only one candidate clears both the theoretical and econometric gates: the capital composition state $\tau_t = \ln K^{\text{ME}} - \ln K^{\text{NRC}}$.

\textit{On theory}: $\tau_t$ indexes the relative capital intensity that the induced innovation mechanism operates through. Rising $\tau_t$ means machinery deepening relative to the structural envelope---exactly the choice-of-technique adjustment the model predicts when wages rise. It does not require labor data ($L_t$), which would introduce an additional source of measurement error and stochastic contamination.

\textit{On integration order}: the cointegration evidence in Tables~\ref{tab:vecm_drift_ledger} and \ref{tab:vecm_trend_ledger} establishes that $k^{\text{ME}}$ and $k^{\text{NRC}}$ cointegrate with a coefficient near unity under drift. A unit cointegrating vector means $\tau_t = k^{\text{ME}} - k^{\text{NRC}}$ is stationary or borderline-stationary---precisely the condition that makes it safe as a regressor in an $I(1)$ cointegration design alongside $k^{\text{NRC}}$. The centered composition variable $\tilde{\tau}_t$ enters specification $D_8$/$D_9$ with manageable condition numbers (4.81--6.31) and passes the polynomial-cointegration gate's residual-stationarity threshold.

The raw mechanization intensity $q_t = K^{\text{ME}}/L$ fails on integration order grounds despite being the direct observable counterpart of the FOC. This is not a defect of the theory. The theory describes the firm's decision mechanism---firms respond to distributional pressure by adjusting $q^*$. The econometric problem is that the \textit{observable} $q_t$ accumulates two $I(1)$ processes ($K^{\text{ME}}$ and $L$) whose ratio inherits $I(2)$ risk. The accumulated-path proxies fail for a different reason: they are nearly collinear with aggregate capital because they cumulate the same underlying growth series. The theory motivating $Q_t^\omega$ is sound---distributional conditioning of accumulation---but the stock variable it produces is empirically indistinguishable from $k_t^{\text{cap}}$.

The surviving specification family takes the form:
\begin{equation*}
y_t = \alpha + \theta \tilde{k}_t^{\text{NRC}} + \psi \tilde{\tau}_t + \phi \tilde{\omega}_t + \lambda \hat{u}_t^{(\tau, \omega)} + \varepsilon_t, \tag{$D_9$}
\end{equation*}
where $\hat{u}_t^{(\tau,\omega)}$ is the FWL residual from regressing the $\tau_t \cdot \omega_t$ interaction on its lower-order terms. The composition response is $\psi + \lambda \tilde{\omega}_t$: distribution conditions the productivity of capital composition. The scale elasticity $\theta$ is the nonresidential-capital envelope parameter. This is not a fallback. It is the only specification that the data admit.

\FloatBarrier
\section{Conclusion}

The stochastic properties of heterogeneous capital are not a technical prerequisite to be checked and set aside. They encode the institutional history of accumulation---and they constrain what specifications are admissible for estimating the transformation elasticity $\theta_t$. The cointegrating coefficient between machinery and structures carries the political economy of the choice of technique: who controls the decision to deepen mechanization, under what distributional conditions, against what structural constraints.

The specification fork is resolved. The capital composition state $\tau_t = \ln K^{\text{ME}} - \ln K^{\text{NRC}}$ is the sole surviving candidate for carrying mechanization intensity into the $\theta_t$ estimation stage. Raw mechanization intensity ($q_t$), accumulated-path proxies ($Q_t^\omega$, $Q_t^q$), and unrestricted component levels are rejected---not for lack of theoretical motivation, but because the integration order evidence imposes binding constraints that no amount of theoretical appeal can override. The LOESS inflections, drift/trend VECM cross-specifications, bootstrap corrections (including Cavaliere et al.\ 2012 wild bootstrap testing), speed-of-adjustment loading vectors, and the consolidated mechanization verdict table establish the empirical foundation for Chapter 2's distribution-conditioned capacity framework.

\end{document}
